\documentclass[10pt,a4paper]{amsart}
\usepackage[margin=19mm]{geometry}
\usepackage{amsmath,amssymb,mathtools}
\usepackage[scr=rsfs,cal=euler]{mathalfa}
\usepackage{xfrac}
\usepackage[unicode,hidelinks,bookmarks=false]{hyperref}
\newcommand{\ie}{i.e.\ }
\newcommand{\cf}{cf.\ }
\newcommand{\resp}{respectively\ }
\newcommand{\Subsection}{\subsection}
\providecommand{\nobreakdash}{\nobreak\hbox{-}}
\setlength{\emergencystretch}{2em}
\multlinegap0pt
\usepackage{bm}
\usepackage{xcolor}
\usepackage{amssymb}
\usepackage[shortlabels]{enumitem}
\usepackage{float}
\usepackage{cancel}
\usepackage[normalem]{ulem}
\usepackage{mathtools}
\usepackage{algorithm}\usepackage{algpseudocode}
\usepackage{tikz}
\usepackage{tcolorbox}
\newtheorem{proposition}{Proposition}
\usepackage{cleveref}
\crefname{proposition}{Proposition}{Propositions}
\def \Fhti {\bm{F}_\htt^i}
\def \Shti {\calS_\htt^i}
\def \calN {\mathcal{N}}
\def \calP {\mathcal{P}}
\def \calQ {\bm{\mathcal{Q}}}
\def \calS {\mathcal{S}}
\def \calT {\mathcal{T}}
\def \f  {\varphi}
\def \flx {\bm{\sigma}_{\htt}^i}
\def \flxa {\bm{\sigma}^a_{\htt}}
\def \htt {{h\tau}}
\def \l  {\lambda}
\def \nhat {\bm{n}_{\mathrm{x}}}
\def \oma {\omega_a}
\def \p  {\partial}
\def \t  {\tau}
\def \x {{\bm{x}}}
\title{An adaptive, space-time discretized linear iterative scheme for doubly-degenerate parabolic problems}
\author{Ayesha Javed, Koondanibha Mitra, Iuliu Sorin Pop}
\date{Source: arXiv:2607.05714v1 (CC BY 4.0)}
\begin{document}
\maketitle

\noindent\emph{Source and licence.} An adaptive, space-time discretized linear iterative scheme for doubly-degenerate parabolic problems. Version arXiv:2607.05714v1 (CC BY 4.0), distributed by arXiv under the Creative Commons Attribution 4.0 International licence, http://creativecommons.org/licenses/by/4.0/, as recorded in the arXiv licence metadata for this exact version and shown on the abstract page https://arxiv.org/abs/2607.05714v1. Full licence notice: \texttt{licenses/arxiv-2607.05714-CC-BY-4.0.txt}.\par\smallskip

\noindent\textit{Editorial scope.} Counted unit: the article's proposition mass\_balance\_property (source0.tex lines 1819--1824). The article's complete original proof of it is delivered with it (source0.tex lines 1825--1869, one \texttt{proof} environment of 45 lines). No mathematics is written, repaired or re-derived: the statement and the proof are the article's own lines, in the article's own notation, and no retained line is substituted or re-flowed.  Retained uncounted as the source-local context the proof needs: lines 1798--1802; lines 1810--1813.  One declared adaptation: exactly 1 ancillary strings inside the retained spans are omitted (image-inclusion commands and, where the source repeats a label, the repeated label definitions) and no other line differs from the source; the figure environments, with their placement, captions and labels, are kept, and the package records the omitted strings in fidelity receipts bound to the affected bodies.  The retained lines cite 1 works, whose entries are transcribed verbatim from the article's own bibliography source in the e-print and delivered with short editorial labels recorded in the item's reference mapping.  The retained lines contain the article's own diagram code, which the wrapper typesets with the standard tikz packages; no image file is delivered and the package declares no asset dependency.  Editorial formatting only, in addition: an AMS \texttt{amsart} preamble that restates the article's own declarations and macros used by the retained lines, namely the environments \texttt{proposition} on separate counters, together with the standard packages those lines need.  The retained environments print in this document as Proposition 1 in order of appearance, whereas the article prints the same items with the numbering of its own full document.  The complete native source of the article as distributed in the arXiv e-print for this exact version is bundled unchanged as \texttt{sources/204567/source0.tex}.
\begin{align}\label{eq:source_flx}
 \int_0^T ( \Shti, \varphi_{h\tau} ) \, dt
 + \int_0^T ( \Fhti, \nabla \varphi_{h\tau} ) \, dt = 0,
 \quad \;\text{for all}\quad\; \varphi_{h\tau} \in \mathcal{V}_{h\tau},
\end{align}

    \begin{align}\label{eq:optimization}
        \flxa:= \underset{\substack{\bm{v}^a_h\in\calQ_\htt^a\\
       \nabla\cdot \bm{v}^a_h = -\Pi_h^0 (\phi_a \Shti) - \nabla\phi_a \cdot \bm{\Pi_h^0} \Fhti}}{\mathrm{argmin}} \left\| \bm{v}^a_h + \phi_a \Fhti\right\|_{L^2(\oma)}
    \end{align}

\begin{proposition}[Well-posedness and mass balance property of $\flx$]\label{mass_balance_property}
    The equilibrated flux $\flx \in \calQ_\htt$ is well defined and satisfies 
    \[
    \nabla\cdot \flx= -\Pi_h^0(\Shti) \qquad \text{ in all } K\in \calT_\htt.
    \]
\end{proposition}

\begin{proof}
First, by counting degrees of freedom for a given $a\in \calN_\htt$ we show that there are indeed $\bm{v}_h^a\in \calQ_\htt^a $ such that $\nabla\cdot \bm{v}^a_h = -\Pi_h^0 (\phi_a \Shti) - \nabla\phi_a \cdot \bm{\Pi_h^0} \Fhti\in \calP_0(\calT_\htt)$. \textcolor{black}{For a given vertex $a \in \mathcal{N}_{h\tau}$ such that $a \notin \Gamma_T$, let $N_a \in \mathbb{N}$ denote the number of elements in $\omega_a$. This gives $N_a$ constraints for satisfying the divergence conditions. Observe that $\partial\omega_a$ contains $N_a$ faces. Since the elements in $\mathcal{T}_{h\tau}$ are $(d+1)$-simplices, each of the remaining $(d+1)N_a$ faces is shared by two adjacent elements, giving in total $((d+1)N_a/2)$ internal faces}. 
%Moreover, $\p\oma$ contains $N_a$ faces since every element has exactly one face not including $a$, and thus lying on the boundary. The remaining $(d+1)N_a$ faces of elements (since the elements are $(d+1)$-dimensional simplices),  are all shared between two elements, thus having $((d+1)N_a/2)$ internal faces. 
For $\bm{v}_h^a\in \calQ_\htt^a$, the trace $\bm{v}^a_h\cdot \nhat$ on each face $E$ gives a polynomial of degree $p$ with $d$ variables (since $t=e-\nhat\cdot \x$ can be substituted as the equation of the face). This implies at most $C_p^{p+d}$ unknown coefficients per  face, thus $((d+1)/2 +1) N_a C_p^{p+d}$ coefficients in total. Now, the total degrees of freedom of $\bm{v}_h^a$ (which is $d$-dimensional) per element are $dC^{p+d+1}_p$. The total number of degrees of freedom needs to exceed the number demanded by constraints, which yields the relation
\[
dC^{p+d+1}_p N_a \geq \left[\left(\frac{d+1}{2} +1\right) C_p^{p+d} + 1\right]N_a,
\]
or, after cancelling terms,
\[
\frac{d(p+d+1)}{d+1}\geq \frac{d+3}{2} + \frac{1}{C_p^{p+d}}.
\]
Rearranging, we get the condition \[p\geq \frac{1}{2}(3-d)(d+1) + \frac{d+1}{d C_p^{p+d}}\]
which is satisfied for all $p\geq 2$ if $d\geq 2$. For $d=1$, this restriction is actually an overcount since $\bm{v}_h^a$ is scalar, and thus $\bm{v}^a_h\cdot \nhat=0$ implies $\bm{v}_h^a=0$. This condition has to be enforced in the boundary nodes of $\oma$, see \Cref{fig:P2nodes}. The value at the internal nodes, $(N_a+1)$ in number, can be freely chosen to satisfy $N_a$ divergence constraints (actually one less due to the compatibility condition). When $a\in \calN_\htt$ is on the boundary $\Gamma_T$, then the restrictions are even less, since no conditions are required on $\p\oma \cap \Gamma$. Thus, for $p\geq 2$ the space $\calQ_\htt^a$  has the required degrees of freedom to satisfy conditions for $\bm{v}_h^a$ in \eqref{eq:optimization}.

    \begin{figure}[h]
                \centering
                    
                    \caption{For one space dimension ($d=1$), the  $P2$ nodal points of $\oma$.}
        \label{fig:P2nodes}
    \end{figure}



Next, we verify the compatibility criteria for $\bm{v}^a_h\in \calQ^a_\htt$. Using the divergence theorem on the field $(\bm{v}^a_h,0)$, we have
\begin{align*}
        0&=\int_{\p\oma} \bm{v}^a_h\cdot \nhat= \int_{\oma}  \nabla \cdot \bm{v}^a_h = - \left[\int_{\oma} \Pi_h^0 (\phi_a \Shti) + \nabla\phi_a \cdot \bm{\Pi_h^0} \Fhti\right]\\
        &= -\left[\int_0^T (\Shti, \phi_a ) \, dt
 + \int_0^T ( \Fhti, \nabla \phi_a ) \, dt\right]=0.
\end{align*}
The last two steps follow from the property of the projector $\Pi^0_h$, and from \eqref{eq:source_flx} by inserting $\f_\htt=\phi_a$. The Euler-Lagrange equation corresponding to the constrained minimization problem \eqref{eq:optimization} then becomes: find $(\flxa,\l^a_\htt)\in \calQ_\htt^a \times \calP_1(\calT_a)$ such that 
\begin{equation*}
    \left\{\begin{aligned}
        &(\flxa + \phi_a \Fhti, \bm{v}^a_h)_{L^2(\oma)} + (\l^a_\htt, \nabla \cdot \bm{v}^a_h)_{L^2(\oma)}= 0,\\
        &(l_h^a, \nabla \cdot \flxa)_{L^2(\oma)}= -(l_h^a\phi_a, \Shti)_{L^2(\oma)}- (l_h^a \nabla \phi_a, \Fhti)_{L^2(\oma)},
    \end{aligned}\right.  
\end{equation*}
for all $(\bm{v}^a_h,l_h^a)\in \calQ_\htt^a \times \calP_1(\calT_a)$. The existence-uniqueness of a solution to this system follows from the standard inf-sup condition since it is a saddle-point problem \cite[Chapter 4]{boffi2013mixed}.

Finally, we get by the partition of unity property $\sum_{a\in \calN_{h\t}} \phi_a=1$ that
\begin{align*}
    &\nabla \cdot \flx = \sum_{a\in \calN_{h\t}} \nabla \cdot \flxa= -\sum_{a\in \calN_{h\t}} \left[\,\Pi_h^0 (\phi_a \Shti) + \nabla\phi_a \cdot \bm{\Pi_h^0} \Fhti\right]\\
    &= - \left[\Pi_h^0 \left((\Sigma_{a\in \calN_{h\t}} \phi_a) \Shti\right) + \nabla(\Sigma_{a\in \calN_\htt}\phi_a ) \cdot \bm{\Pi_h^0} \Fhti\right]= - \Pi_h^0 \Shti.
\end{align*}
%This concludes the proof.
\end{proof}

\begin{thebibliography}{99}
\bibitem[boffi2]{boffi2013mixed} D. Boffi, F. Brezzi, and M. Fortin, \emph{{Mixed Finite Element Methods} and {Applications}}, \textbf{44}, (2013), Springer. \href{https://doi.org/10.1007/978-3-642-36519-5}{DOI:10.1007/978-3-642-36519-5}
\end{thebibliography}
\end{document}
